\section{总结与延伸}

\subsection{核心要点回顾}

本次演讲系统地梳理了生成式3D方法的演进脉络和当前格局：

\begin{enumerate}
    \item \textbf{3D生成方法的演进}：从需要3D真值的3D GAN，到利用可微渲染摆脱3D数据依赖的 PlatonicGAN，再到多视角融合的 PixelNeRF
    \item \textbf{视频模型的3D理解}：实验证实视频模型能隐式建模3D几何，但推理成本限制了部署
    \item \textbf{CAT3D}：「多视角扩散 + NeRF 蒸馏」范式，约1分钟生成3D场景
    \item \textbf{B3D}：前馈 Gaussian Head 替代 NeRF 优化，7秒内生成3D场景
    \item \textbf{几何 VAE}：图像 VAE 不适用于几何数据，需要专门的架构和损失函数设计
    \item \textbf{VLM 感知对齐}：用 VLM 作为评判者进行 DPO 训练，改善自编码器的感知质量
    \item \textbf{ROAR}：将重光照作为额外控制维度引入3D生成框架
    \item \textbf{Captive}：立方体贴图 + 冻结基础模型实现数据高效的全景生成
\end{enumerate}

\subsection{关键设计原则}

\begin{importantbox}{3D生成方法的设计原则}
\textbf{1. 2D生成 + 3D蒸馏}：利用2D模型的强大生成能力，通过3D重建获得高效可渲染的3D表示。

\textbf{2. 分阶段设计}：将容易获得监督的属性（如 RGB、XYZ）用扩散模型生成，将难以获得真值的属性（如 opacity、scale）用前馈网络回归。

\textbf{3. 冻结 + 适配}：在基础模型上冻结已学习的强大先验，仅训练少量适配层来学习新约束。

\textbf{4. 跨视角注意力}：通过全局注意力机制确保多视角生成的3D一致性。
\end{importantbox}

\subsection{拓展阅读}

\begin{itemize}
    \item CAT3D: Create Anything in 3D with Multi-View Diffusion Models (Gao et al., Google, 2024)
    \item B3D: 快速前馈3D场景生成方法 (Google Research)
    \item Splatter Image: Ultra-Fast Single-View 3D Reconstruction (Szymanowicz et al., 2024)
    \item MASt3R: 3D Structure-from-Motion 管线
    \item ROAR: 3D重光照方法 (Google Research)
    \item The Bitter Lesson --- Richard Sutton, 2019
    \item Genie 2: 交互式视频生成模型 (Google DeepMind)
    \item World Labs: 实时3D世界模型
\end{itemize}

%% --- Fin del contenido --- %%

\end{document}
